Autonomous Surface Vehicles (ASVs) are increasingly deployed in unstructured maritime environments such as busy ports and marinas. These domains present unique perception challenges, including dynamic water reflections, waves, and low-profile hazards. %such as floating debris. 
While high-end platforms rely on active sensors like marine radar or LiDAR, small-scale ASVs demand low-cost, vision-only solutions that operate reliably under strict power and hardware constraints.

In safety-critical navigation, a gap exists between perception accuracy and operational safety: high accuracy in perception modules, such as object detection and segmentation, does not inherently translate into safe navigation. The fundamental question for an ASV is not whether an obstacle is perfectly segmented, but a more urgent one: \emph{``Does this obstacle pose an immediate collision risk?"} A vision system that returns raw detections without this risk-aware context can overwhelm the controller with data while failing to highlight the most pressing threats, such as treating a distant, harmless buoy and a swimmer in the immediate path as equally salient.

To bridge this gap, we propose WARD (Water-surface Awareness and Risk Detection), a vision-only (RGB), risk-aware framework that fuses SegFormer~\cite{xie2021segformersimpleefficientdesign} semantic segmentation with radar-calibrated Depth Anything V2~\cite{yang2024depthv2} monocular depth, as illustrated in Fig.~\ref{fig:pipeline}. WARD introduces a continuous collision-corridor risk model to reason jointly about proximity, lateral offset, closing velocity, and semantic hazard. By delivering actionable, real-time safety assessments (SAFE, CAUTION, and IMMEDIATE) from a single RGB stream, WARD offers a practical baseline for resource-constrained ASVs operating in dynamic environments.

% In summary, this paper contributes:
% \begin{itemize}
%     \item an end-to-end, vision-only pipeline that fuses semantic segmentation and monocular depth through connected-component instance extraction, avoiding specialized ranging hardware;
%     \item a radar-calibrated scaling strategy that restores near-range metric accuracy in a depth model pre-trained on road-driving imagery, within the band that matters most for collision warning;
%     \item a continuous collision-corridor risk score that unifies RSS-style geometric reasoning with semantic class weighting and closing-velocity amplification, in place of a binary safe/unsafe decision; and
%     \item a depth-augmented SORT tracker that jointly estimates bounding-box motion and closing velocity, yielding stable Time-to-Collision estimates despite noisy per-frame depth.
% \end{itemize}

%
% We evaluate WARD on real-world maritime sequences, prioritizing safety-oriented metrics over segmentation accuracy alone and ensuring that the system remains reliable and viable for real-time deployment on resource-constrained ASVs.